\section{Experiments}

\subsection{Experimental Setup}
\textbf{Dataset:} A curated dataset of 149 synthesis-property relationships across 85 materials systems (in-domain and out-of-domain).
\textbf{Baselines:} Baseline LLM (Gemini-1.5-Pro), Naive KG+LLM (Standard RAG), Online KG+LLM.
\textbf{Evaluation:} We employ a \textbf{Constraint-Based Rubric} using an LLM-judge. Unlike open-ended evaluation, this judge verifies specific physical constraints (e.g., "Does the proposed temperature fall within the 700-900C range?"), which has been shown to correlate highly with human expert judgment \cite{zheng_chatgpt_2023}.

\subsection{Results & Discussion}

\begin{figure}[h]
    \centering
    \includegraphics[width=\linewidth]{iclr2026/Fig/ICLR-Fig4-2.pdf} 
    \caption{\textbf{Performance comparison.} \texttt{ARIA} recovers performance lost by Naive KG integration.}
    \label{fig:results}
\end{figure}

\subsubsection{Contextual Tunneling is Systemic}
As shown in Table \ref{tab:main_results}, \textbf{Naive KG+LLM consistently underperforms the Baseline LLM} (e.g., -28.7% in inverse design). This confirms that "more knowledge" can be detrimental if not integrated causally. The model fixates on retrieved similarities (Tunneling) rather than reasoning.

\subsubsection{ARIA's Rescue Effect}
\texttt{ARIA} successfully reverses this degradation. In complex inverse design tasks, it improves upon Naive KG+LLM by \textbf{36.6%}, restoring performance to levels comparable to the Baseline.

\subsubsection{The "Why RAG?" Paradox}
A key observation is that the Baseline LLM often performs best numerically. Why use RAG?
\begin{itemize}
    \item \textbf{Black Box vs. Glass Box:} The Baseline provides no provenance. In scientific discovery, a hallucinated "correct" answer is dangerous.
    \item \textbf{Provenance:} \texttt{ARIA} provides verifiable graph paths (Tier 1) and analogical reasoning chains (Tier 2).
    \item \textbf{Conclusion:} \texttt{ARIA} achieves the \textit{reliability} of RAG without the \textit{performance penalty} of Contextual Tunneling. It offers a "safe" path to using KGs in science.
\end{itemize}

\begin{table}[h]
\caption{Overall Performance Comparison (Inverse Design)}
\label{tab:main_results}
\begin{tabular}{l|ccc}
\toprule
Method & Accuracy & Reasoning & Overall \\
\midrule
Baseline LLM & \textbf{0.64} & \textbf{0.64} & \textbf{0.59} \\
Naive KG+LLM & 0.47 & 0.41 & 0.42 \\
Online KG+LLM & 0.51 & 0.48 & 0.45 \\
\textbf{ARIA (Ours)} & 0.63 & 0.59 & 0.57 \\
\bottomrule
\end{tabular}
\end{table}

\todo{Update table values with specific data from the `table2.tex` file if available, or use placeholders if exact numbers are needed.}
